\documentclass[10pt]{article}\usepackage[]{graphicx}\usepackage[]{xcolor}
% maxwidth is the original width if it is less than linewidth
% otherwise use linewidth (to make sure the graphics do not exceed the margin)
\makeatletter
\def\maxwidth{ %
  \ifdim\Gin@nat@width>\linewidth
    \linewidth
  \else
    \Gin@nat@width
  \fi
}
\makeatother

\definecolor{fgcolor}{rgb}{0.345, 0.345, 0.345}
\newcommand{\hlnum}[1]{\textcolor[rgb]{0.686,0.059,0.569}{#1}}%
\newcommand{\hlsng}[1]{\textcolor[rgb]{0.192,0.494,0.8}{#1}}%
\newcommand{\hlcom}[1]{\textcolor[rgb]{0.678,0.584,0.686}{\textit{#1}}}%
\newcommand{\hlopt}[1]{\textcolor[rgb]{0,0,0}{#1}}%
\newcommand{\hldef}[1]{\textcolor[rgb]{0.345,0.345,0.345}{#1}}%
\newcommand{\hlkwa}[1]{\textcolor[rgb]{0.161,0.373,0.58}{\textbf{#1}}}%
\newcommand{\hlkwb}[1]{\textcolor[rgb]{0.69,0.353,0.396}{#1}}%
\newcommand{\hlkwc}[1]{\textcolor[rgb]{0.333,0.667,0.333}{#1}}%
\newcommand{\hlkwd}[1]{\textcolor[rgb]{0.737,0.353,0.396}{\textbf{#1}}}%
\let\hlipl\hlkwb

\usepackage{framed}
\makeatletter
\newenvironment{kframe}{%
 \def\at@end@of@kframe{}%
 \ifinner\ifhmode%
  \def\at@end@of@kframe{\end{minipage}}%
  \begin{minipage}{\columnwidth}%
 \fi\fi%
 \def\FrameCommand##1{\hskip\@totalleftmargin \hskip-\fboxsep
 \colorbox{shadecolor}{##1}\hskip-\fboxsep
     % There is no \\@totalrightmargin, so:
     \hskip-\linewidth \hskip-\@totalleftmargin \hskip\columnwidth}%
 \MakeFramed {\advance\hsize-\width
   \@totalleftmargin\z@ \linewidth\hsize
   \@setminipage}}%
 {\par\unskip\endMakeFramed%
 \at@end@of@kframe}
\makeatother

\definecolor{shadecolor}{rgb}{.97, .97, .97}
\definecolor{messagecolor}{rgb}{0, 0, 0}
\definecolor{warningcolor}{rgb}{1, 0, 1}
\definecolor{errorcolor}{rgb}{1, 0, 0}
\newenvironment{knitrout}{}{} % an empty environment to be redefined in TeX

\usepackage{alltt}

\usepackage[T1]{fontenc}
\usepackage[utf8]{inputenc}
\usepackage{amsmath,amssymb}
\usepackage{array,booktabs,float}
\usepackage[margin=1in]{geometry}
\usepackage[hidelinks]{hyperref}

\setlength{\parindent}{0pt}
\setlength{\parskip}{0.65em}
\setlength{\emergencystretch}{2em}

\title{Variable Selection and MNARz Clustering\\
with \texttt{SelvarMixMNAR}}
\author{}
\date{}

%\VignetteIndexEntry{Variable Selection and MNARz Clustering with SelvarMixMNAR}
%\VignetteEngine{knitr::knitr}
%\VignetteEncoding{UTF-8}
\IfFileExists{upquote.sty}{\usepackage{upquote}}{}
\begin{document}

\maketitle



\texttt{SelvarMixMNAR} is the direct R implementation accompanying the
NeurIPS 2025 paper \emph{A Unified Framework for Variable Selection in
Model-Based Clustering with Missing Not at Random}
\cite{ho2025unifiedmbcmnar}. The package combines a penalized Gaussian-mixture
ranking, SRUW variable roles, and a MNARz missingness model.

For incomplete data, \texttt{decoupled\_mnarz} estimates SRUW roles from a
completed working matrix before fitting an MNARz mixture, whereas
\texttt{joint\_mnarz} estimates the SRUW and missingness terms under one
observed-data criterion.

\section{SRUW roles and the fitted MNARz mechanism}

For observation $Y_i \in \mathbb{R}^p$, the SRUW factorization uses three
disjoint role sets and one predictor subset:

\begin{itemize}
\item $S$: variables that define the Gaussian mixture;
\item $R \subseteq S$: clustering variables used as predictors of redundant
 variables;
\item $U$: redundant variables modeled by a common Gaussian regression on
 $Y_i^R$;
\item $W$: variables modeled independently of the clustering block.
\end{itemize}

Writing $K$ for the number of mixture components, the factorization has the
form
\[
f(Y_i)
= f_{\mathrm{clust}}(Y_i^S \mid K,m)
 f_{\mathrm{reg}}(Y_i^U \mid Y_i^R,r)
 f_{\mathrm{indep}}(Y_i^W \mid l).
\]
Here $m$, $r$, and $l$ index the covariance models for the clustering,
regression, and independent blocks, respectively.

The fitted missingness class is MNARz. If $C_{ij}=1$ denotes a missing entry
and $Z_i=k$ denotes latent component membership, then
\[
\Pr(C_{ij}=1 \mid Z_i=k)=\rho_k.
\]
Thus each component has one missingness probability shared across variables.
MNARy, in which missingness depends directly on an unobserved value, and
MNARzj, in which the class effect varies by coordinate, require different
likelihoods. Identifiability relies on the MNARz assumptions stated in the
accompanying paper in addition to numerical convergence.

The separate decoupled fit also permits a prespecified mixed mask through
\texttt{mnarz\_control = list(mecha = "mixed", is\_mnar = m)}. All
coordinates with $m_j=\mathrm{TRUE}$ share the class-specific parameter
$\rho_k$. Mask factors with $m_j=\mathrm{FALSE}$ are treated as ignorable and
are omitted from the modeled likelihood. The vector $m$ is supplied by the
analyst, and all selected coordinates share $\rho_k$. It is therefore a
 sensitivity specification rather than MNARzj.

\section{Two stages of variable selection}

The ranking stage fits penalized Gaussian mixtures over a grid of mean and
precision penalties. A variable receives a high score when its component
means remain nonzero across many valid grid fits. Each candidate value of $K$
owns its initializer, penalty grid, fitted states, and failure records.

The role-selection stage scans the ranked variables. A variable enters $S$
when its clustering contribution improves the model criterion. Variables not
retained in $S$ are assigned to $U$ or $W$ by comparing common regression and
independent Gaussian models. The stopping parameter \texttt{hsize} is the
number of consecutive nonpositive criterion increments tolerated before the
scan stops; an improvement resets the count.

\section{A complete-data example}

The following Gaussian example has two clustering variables, two redundant
variables, and two independent variables. The generating labels are retained
only for external evaluation and are not supplied to the estimator.

\begin{knitrout}
\definecolor{shadecolor}{rgb}{0.969, 0.969, 0.969}\color{fgcolor}\begin{kframe}
\begin{alltt}
\hlkwd{library}\hldef{(SelvarMixMNAR)}
\end{alltt}


{\ttfamily\noindent\bfseries\color{errorcolor}{\#> Error in `library()`:\\\#> ! there is no package called 'SelvarMixMNAR'}}\begin{alltt}
\hlkwd{set.seed}\hldef{(}\hlnum{2025}\hldef{)}
\hldef{n} \hlkwb{<-} \hlnum{90}
\hldef{latent_class} \hlkwb{<-} \hlkwd{rep}\hldef{(}\hlnum{1}\hlopt{:}\hlnum{3}\hldef{,} \hlkwc{each} \hldef{= n} \hlopt{/} \hlnum{3}\hldef{)}
\hldef{signal_1} \hlkwb{<-} \hlkwd{rnorm}\hldef{(n,} \hlkwd{c}\hldef{(}\hlopt{-}\hlnum{2}\hldef{,} \hlnum{0}\hldef{,} \hlnum{2}\hldef{)[latent_class],} \hlnum{0.7}\hldef{)}
\hldef{signal_2} \hlkwb{<-} \hlkwd{rnorm}\hldef{(n,} \hlkwd{c}\hldef{(}\hlnum{1.5}\hldef{,} \hlopt{-}\hlnum{1.5}\hldef{,} \hlnum{0}\hldef{)[latent_class],} \hlnum{0.7}\hldef{)}
\hldef{x} \hlkwb{<-} \hlkwd{cbind}\hldef{(}
 \hlkwc{signal_1} \hldef{= signal_1,}
 \hlkwc{signal_2} \hldef{= signal_2,}
 \hlkwc{redundant_1} \hldef{=} \hlnum{0.8} \hlopt{*} \hldef{signal_1} \hlopt{+} \hlkwd{rnorm}\hldef{(n,} \hlkwc{sd} \hldef{=} \hlnum{0.35}\hldef{),}
 \hlkwc{redundant_2} \hldef{=} \hlopt{-}\hlnum{0.7} \hlopt{*} \hldef{signal_2} \hlopt{+} \hlkwd{rnorm}\hldef{(n,} \hlkwc{sd} \hldef{=} \hlnum{0.35}\hldef{),}
 \hlkwc{noise_1} \hldef{=} \hlkwd{rnorm}\hldef{(n),}
 \hlkwc{noise_2} \hldef{=} \hlkwd{rnorm}\hldef{(n)}
\hldef{)}
\hlkwd{dim}\hldef{(x)}
\end{alltt}
\begin{verbatim}
#> [1] 90  6
\end{verbatim}
\end{kframe}
\end{knitrout}

The known variable order is supplied through \texttt{rank} to isolate the SRUW
role-selection stage in this synthetic example. Omitting \texttt{rank}
estimates the penalized ordering from the data.

\begin{knitrout}
\definecolor{shadecolor}{rgb}{0.969, 0.969, 0.969}\color{fgcolor}\begin{kframe}
\begin{alltt}
\hlkwd{set.seed}\hldef{(}\hlnum{2026}\hldef{)}
\hldef{fit} \hlkwb{<-} \hlkwd{SelvarClustLasso}\hldef{(}
 \hlkwc{x} \hldef{= x,}
 \hlkwc{nbcluster} \hldef{=} \hlnum{3}\hldef{,}
 \hlkwc{models} \hldef{=} \hlsng{"VVI"}\hldef{,}
 \hlkwc{rank} \hldef{=} \hlkwd{seq_len}\hldef{(}\hlkwd{ncol}\hldef{(x)),}
 \hlkwc{hsize} \hldef{=} \hlnum{2}\hldef{,}
 \hlkwc{nbcores} \hldef{=} \hlnum{1}\hldef{,}
 \hlkwc{workflow} \hldef{=} \hlsng{"imputed_sruw"}\hldef{,}
 \hlkwc{init_control} \hldef{=} \hlkwd{list}\hldef{(}\hlkwc{method} \hldef{=} \hlsng{"mclust_hc"}\hldef{,} \hlkwc{seed} \hldef{=} \hlnum{2026}\hldef{),}
 \hlkwc{selection_control} \hldef{=} \hlkwd{list}\hldef{(}\hlkwc{stopping} \hldef{=} \hlsng{"consecutive"}\hldef{,} \hlkwc{seed} \hldef{=} \hlnum{2026}\hldef{)}
\hldef{)}
\end{alltt}


{\ttfamily\noindent\bfseries\color{errorcolor}{\#> Error in `SelvarClustLasso()`:\\\#> ! could not find function "{}SelvarClustLasso"{}}}\end{kframe}
\end{knitrout}

The \texttt{imputed\_sruw} workflow fits the ranking and SRUW role model on the
complete working data without a missingness likelihood. The
default Mclust covariance model is \texttt{VVI}, a diagonal model with
component-specific volume and shape.

\section{Decoupled and joint MNARz workflows}

When \texttt{x} is incomplete, \texttt{workflow = "auto"} resolves to
\texttt{"decoupled\_mnarz"}. Table~\ref{tab:workflows} separates the three
statistical targets.

\begin{table}[H]
\centering
\small
\caption{Workflow-specific statistical targets.}
\label{tab:workflows}
\begin{tabular}{@{}>{\raggedright\arraybackslash}p{0.24\linewidth}
          >{\raggedright\arraybackslash}p{0.68\linewidth}@{}}
\toprule
Workflow & Target \\
\midrule
\texttt{imputed\_sruw} &
SRUW ranking and role selection on a completed working data set; no
missingness likelihood. \\
\texttt{decoupled\_mnarz} &
Imputed-data SRUW selection followed by a separate, role-agnostic MNARz
Gaussian mixture. \\
\texttt{joint\_mnarz} &
Observed-data SRUW and MNARz estimation on the original incomplete matrix,
with forward and reverse role updates scored by total observed-data BIC. \\
\bottomrule
\end{tabular}
\end{table}

The decoupled estimator reduces computation by separating imputed-data role
selection from the missingness fit. The joint estimator couples both terms in
the observed-data objective and consequently requires more repeated mixture
fits along the role path.

\begin{knitrout}
\definecolor{shadecolor}{rgb}{0.969, 0.969, 0.969}\color{fgcolor}\begin{kframe}
\begin{alltt}
\hlkwd{set.seed}\hldef{(}\hlnum{2027}\hldef{)}
\hldef{x_miss} \hlkwb{<-} \hldef{x}
\hldef{rho} \hlkwb{<-} \hlkwd{c}\hldef{(}\hlnum{0.02}\hldef{,} \hlnum{0.08}\hldef{,} \hlnum{0.15}\hldef{)}
\hldef{mask_probability} \hlkwb{<-} \hlkwd{matrix}\hldef{(}
 \hldef{rho[latent_class],} \hlkwc{nrow} \hldef{=} \hlkwd{nrow}\hldef{(x_miss),} \hlkwc{ncol} \hldef{=} \hlkwd{ncol}\hldef{(x_miss)}
\hldef{)}
\hldef{mask} \hlkwb{<-} \hlkwd{matrix}\hldef{(}\hlkwd{runif}\hldef{(}\hlkwd{length}\hldef{(x_miss)),} \hlkwc{nrow} \hldef{=} \hlkwd{nrow}\hldef{(x_miss))} \hlopt{<} \hldef{mask_probability}
\hldef{x_miss[mask]} \hlkwb{<-} \hlnum{NA_real_}

\hldef{common} \hlkwb{<-} \hlkwd{list}\hldef{(}
 \hlkwc{nbcluster} \hldef{=} \hlnum{2}\hlopt{:}\hlnum{4}\hldef{,}
 \hlkwc{models} \hldef{=} \hlsng{"VVI"}\hldef{,}
 \hlkwc{num_vals_penalty} \hldef{=} \hlnum{3}\hldef{,}
 \hlkwc{hsize} \hldef{=} \hlnum{2}\hldef{,}
 \hlkwc{nbcores} \hldef{=} \hlnum{1}\hldef{,}
 \hlkwc{init_control} \hldef{=} \hlkwd{list}\hldef{(}\hlkwc{method} \hldef{=} \hlsng{"mclust_hc"}\hldef{,} \hlkwc{seed} \hldef{=} \hlnum{2027}\hldef{),}
 \hlkwc{selection_control} \hldef{=} \hlkwd{list}\hldef{(}\hlkwc{stopping} \hldef{=} \hlsng{"consecutive"}\hldef{,} \hlkwc{seed} \hldef{=} \hlnum{2027}\hldef{)}
\hldef{)}

\hldef{fit_decoupled} \hlkwb{<-} \hlkwd{do.call}\hldef{(}
 \hldef{SelvarClustLasso,}
 \hlkwd{c}\hldef{(}\hlkwd{list}\hldef{(}\hlkwc{x} \hldef{= x_miss,} \hlkwc{workflow} \hldef{=} \hlsng{"decoupled_mnarz"}\hldef{), common)}
\hldef{)}

\hldef{fit_joint} \hlkwb{<-} \hlkwd{do.call}\hldef{(}
 \hldef{SelvarClustLasso,}
 \hlkwd{c}\hldef{(}\hlkwd{list}\hldef{(}\hlkwc{x} \hldef{= x_miss,} \hlkwc{workflow} \hldef{=} \hlsng{"joint_mnarz"}\hldef{), common)}
\hldef{)}
\end{alltt}
\end{kframe}
\end{knitrout}

Ranking requires one complete working data set. With
\texttt{use\_copula = TRUE}, the package uses \texttt{gcimputeR}; setting
\texttt{use\_copula = FALSE} requests \texttt{missRanger}. The imputation
supplies the ranking stage in both MNARz workflows. The joint fit itself
returns to the original incomplete matrix.

\section{Initialization and warm paths}

Initialization is data-driven and performed separately for each candidate
component count. The control \path{init_control$method} accepts the following
values:
\begin{itemize}
\item \path{mclust_hc} and \path{mclust_em};
\item \path{kmeans} and \path{random};
\item \path{user} and \path{previous_fit};
\item \path{deterministic_multistart}.
\end{itemize}
Every route produces the same numeric state: component proportions,
means, covariances, precisions, responsibilities, partition, objective,
method, seed, and failure status.

Independent cold starts are the ranking default. Warm starts are requested
through \path{rank_control$warm_start}:

\begin{knitrout}
\definecolor{shadecolor}{rgb}{0.969, 0.969, 0.969}\color{fgcolor}\begin{kframe}
\begin{alltt}
\hldef{fit_warm} \hlkwb{<-} \hlkwd{SelvarClustLasso}\hldef{(}
 \hldef{x,}
 \hlkwc{nbcluster} \hldef{=} \hlnum{2}\hlopt{:}\hlnum{4}\hldef{,}
 \hlkwc{workflow} \hldef{=} \hlsng{"imputed_sruw"}\hldef{,}
 \hlkwc{nbcores} \hldef{=} \hlnum{1}\hldef{,}
 \hlkwc{rank_control} \hldef{=} \hlkwd{list}\hldef{(}\hlkwc{warm_start} \hldef{=} \hlsng{"both"}\hldef{)}
\hldef{)}
\end{alltt}
\end{kframe}
\end{knitrout}

An inner warm start continues the graphical-lasso solver between EM
iterations. An outer warm start passes the complete fitted mixture state
between adjacent points of a deterministic penalty path. Points within one
outer path are serial because each depends on its predecessor; parallelism is
reserved for independent component counts or initializations.

\section{Results and diagnostics}

A fitted \texttt{selvarmix} object provides standard methods:

\begin{knitrout}
\definecolor{shadecolor}{rgb}{0.969, 0.969, 0.969}\color{fgcolor}\begin{kframe}
\begin{alltt}
\hldef{fit}
\end{alltt}


{\ttfamily\noindent\bfseries\color{errorcolor}{\#> Error:\\\#> ! object 'fit' not found}}\begin{alltt}
\hlkwd{summary}\hldef{(fit)}
\end{alltt}


{\ttfamily\noindent\bfseries\color{errorcolor}{\#> Error:\\\#> ! object 'fit' not found}}\begin{alltt}
\hlkwd{plot}\hldef{(fit)}
\end{alltt}


{\ttfamily\noindent\bfseries\color{errorcolor}{\#> Error:\\\#> ! object 'fit' not found}}\end{kframe}
\end{knitrout}

The default plot displays the counts in the $S$/$U$/$W$ role partition; $R$
is reported separately because it is a subset of $S$. The result also
contains the selected component count, criterion metadata, partition,
posterior probabilities, Gaussian and regression parameters, preprocessing
transformations, and the finite completed working-data matrix. Workflow-specific
fields preserve the MNARz or joint fit, initialization provenance, penalty
grids, convergence status, and role-selection stopping diagnostics.

Failed ranking fits remain visible in the recorded grid state. Fits with
non-finite objectives, invalid component masses, non-positive-definite
covariances, inconsistent covariance and precision matrices,
objective decreases, or unsuccessful graphical-lasso termination are not
scored. This rule prevents numerical failure from appearing as evidence that
a variable has been shrunk to zero.

\section{Remarks}

\begin{enumerate}
\item The MNARz mechanism assumes one component-specific missingness
 probability across coordinates. A scientific application should justify
 this restriction from its data-collection process.
\item The optional mixed mask is fixed by the analyst for the separate
 decoupled fit, and estimation conditions on that partition of the
 coordinates.
\item \texttt{joint\_mnarz} requires more computation than
 \texttt{decoupled\_mnarz}. Comparisons of clustering, selection, and
 imputation accuracy should use replicated designs matched to the scientific
 setting.
\end{enumerate}

\bibliographystyle{plain}
\bibliography{references}

\end{document}
